\section{超对称规范理论}

% [Tutor Note: 原文 PDF 第 22 页（印刷页 21）。本节从头建立 $N=1$ 超对称的旋量约定与超场形式，阅读时宜对照原文 PDF 第 22--23 页式 (49) 前后的约定表；文中所引 [22,23,24] 即 \cite{WB,West,Sohnius}（Wess--Bagger、West、Sohnius 三本/篇标准文献）。]

\begin{keypoints}
\begin{itemize}
  \item \textbf{物理图像：} $N=1$ 超对称把玻色子与费米子组织进同一超多重态（标量多重态、矢量多重态）；超代数表示分为长、短多重态，$M=\sqrt2|Z|$ 的 BPS 态只保持一半超对称。
  \item \textbf{数学工具：} 二分量旋量与 $\epsilon$ 张量约定、$\sigma^{\mu\nu}$、超空间 $(x^\mu,\theta^\alpha,\bar\theta_{\dot\alpha})$ 与超场展开、超协变导数 $D_\alpha,\bar D_{\dot\alpha}$、手征超场 $\Phi$ 与矢量超场 $V$、Wess--Zumino 规范。
  \item \textbf{关键结论：} 无中心荷代数 (\ref{susy-noC}) 与含中心荷代数 (\ref{susy-C})--(\ref{susy-ab})，由此得质量界 $M\geq\sqrt2|Z|$；手征超场展开 (\ref{csfy})、矢量超场分量 (\ref{W-abelian})--(\ref{W-nonabelian})；$W^\alpha W_\alpha|_{\theta\theta}$ (\ref{WW-nonabelian}) 给出规范场动能与 $\theta$ 项。
  \item \textbf{伏笔：} 完整 $N=1$ 拉格朗日量 (\ref{fullSF})/(\ref{fullC}) 消去 $F,D^a$ 辅助场后得标量势 (\ref{Spot})；$N=1$ 超对称不固定标量部分与杨--米尔斯部分的相对归一化，这是下一节 $N=2$ 构造的出发点。
\end{itemize}
\end{keypoints}

本节将解释超对称与超对称场论中与 Witten 和 Seiberg 工作相关的若干方面。我们先说明所用的约定，然后简要描述含中心荷与不含中心荷的超对称代数的表示；随后以量子场的语言讨论 $N=1$ 超对称的表示，并构造具有 $N=1$ 与 $N=2$ 超对称的拉格朗日量。本节大部分内容如今已是标准材料，可参见 \cite{WB,West,Sohnius}。在本节末尾，我们将显式算出含物质与不含物质的 $N=2$ 理论中的中心荷。

\subsection{约定}

% [Tutor Note: 原文 PDF 第 22 页（印刷页 21）首句 "We start by descrbing our conventions" 中 "descrbing" 为原文拼写错误，应为 "describing"。译文按正确词义译出，原文误拼在此注明，未改动原文。]

我们先说明所用的约定。采用平直度规 $\eta_{ab} = {\rm diag}~ (1,-1,-1,-1)$。洛伦兹群 $SL(2, C)\sim SU(2)_L\times SU(2)_R$ 的旋量以带点与不带点两类分量书写，在 $SL(2,C)$ 下按如下方式变换：
$$
\psi'_\alpha = M_{\alpha}^{~\beta} \psi_\beta\,, \qquad
\bar\psi'_{\dot{\alpha}} = M^{*~\dot{\beta}}_{\dot{\alpha}}
\bar{\psi}_{\dot{\beta}}\,.
$$
旋量指标用 $\epsilon$ 张量升降，
$$
\epsilon^{\alpha \beta} = \epsilon^{\dot{\alpha}\dot{\beta}} =
\left( \begin{array}{lr}
0 & 1 \\
-1 & 0
\end{array} \right) =
(i \sigma_{2})\,.
$$
按定义，该张量在 $SL(2,C)$ 变换下不变：$\epsilon^{\alpha \beta} = M_{~\gamma}^{\alpha}
\epsilon^{\gamma \delta}M_{\delta}^{~\beta}$。这可以写成
$M^T\sigma_2 M=\sigma_2$，由此可得 $\sigma_2 M=(M^T)^{-1}\sigma_2$。利用这一点，可把带升指标的旋量的变换写成
$$
\psi'^\alpha =\psi^\beta (M^{-1})^{~\alpha}_{\beta}\,,\qquad
\bar{\psi}'^{\dot{\alpha}} = \bar{\psi}^{\dot{\beta}}
(M^*)^{-1~\dot{\alpha}}_{\dot{\beta}}\,.
$$

% [Tutor Note: 原文 PDF 第 22 页，式 (49) 之前。$M^T\sigma_2M=\sigma_2$ 是由 $\epsilon^{\alpha\beta}=M^\alpha_{~\gamma}\epsilon^{\gamma\delta}M^\delta_{~\beta}$ 与 $\sigma_2=i\epsilon$ 对照得到的矩阵形式；注意 $M_\alpha^{~\beta}$（行 $\alpha$、列 $\beta$）与转置 $(M^T)_\alpha^{~\beta}=M^\beta_{~\alpha}$ 的指标对应，以及 $(M^{-1})^\alpha_{~\beta}$ 的指标位置，稍有不慎就会把升降指标的变换写反。]

现在定义
$$
(\sigma^\mu)_{\alpha \dot{\alpha}} \equiv (1, \vec{\sigma})\,,
$$
于是
$$
\sigma^\mu P_\mu = \left ( \begin{array}{lr}
P^{0} - P^{3} & -P^{1} + i P^{2} \\
-P^{1} - iP^{2} & P^{0} + P^{3}
\end{array} \right )\,, 
$$
且 $\det(\sigma^\mu P_\mu)=P_\mu P^\mu$。我们可以用 $\epsilon$ 张量升高 $\sigma^\mu$ 的指标，并把 $\bar\sigma$ 定义为
$$
(\bar{\sigma}^\mu)^{\dot{\alpha}\alpha}=-(\sigma^\mu)^{\alpha\dot{\alpha}}
=\epsilon^{\dot{\alpha}\dot{\beta}}\epsilon^{\alpha \beta}
(\sigma^\mu)_{\beta \dot{\beta}}\,. 
$$
在数值上，这给出
$$
(\bar{\sigma}^\mu) = (i\sigma_2) (\sigma^\mu)^T (i\sigma_2)^T =
\sigma_2 (\sigma^\mu)^T\sigma_2 = (1, - \vec{\sigma})\,.
$$

% [Tutor Note: 原文 PDF 第 23 页（印刷页 22），$\bar\sigma^\mu$ 的定义式。由 $\epsilon$ 张量两次升指标得 $(\bar\sigma^\mu)^{\dot\alpha\alpha}=-(\sigma^\mu)^{\alpha\dot\alpha}=(i\sigma_2)(\sigma^\mu)^T(i\sigma_2)^T$，数值上为 $(1,-\vec\sigma)$，即 $\bar\sigma^0=\sigma^0=1$、$\bar\sigma^i=-\sigma^i$。这一负号是 (+---) 号差下 $\sigma$ 与 $\bar\sigma$ 关系的关键，后文 $W^\alpha W_\alpha$ 中 $F\widetilde F$ 项的符号由此决定；切勿套用 $(-,+,+,+)$ 文献的约定。]

在这样的约定下，洛伦兹变换由下式生成
\begin{eqnarray}
(\sigma^{\mu\nu})_\alpha^{~\beta} & = & \frac{1}{4}
[\sigma^\mu_{\alpha \dot{\beta}} \bar{\sigma}^{\nu\dot{\beta}\beta}- 
(\mu \leftrightarrow \nu)]\,, \nonumber \\
(\bar{\sigma}^{\mu\nu})^{\dot{\alpha}}_{\dot{~\beta}}&=&\frac{1}{4}
[\bar{\sigma}^{\mu\dot{\alpha}\beta}\sigma^\nu_{\beta\dot{\beta}}-
(\mu\leftrightarrow \nu)]\,.\nonumber
\end{eqnarray}
对旋量的标量积，我们采用如下约定
\begin{eqnarray}
\psi \chi &= & \psi^{\alpha} \chi_{\alpha} = - \psi_{\alpha}
\chi^{\alpha} = \chi^{\alpha} \psi_{\alpha} = \chi \psi\,,\nonumber\\ 
\bar{\psi} \bar{\chi} &=& \bar{\psi}_{\dot{\alpha}} 
\bar{\chi}^{\dot{\alpha}} =\bar{\chi} \bar{\psi}\,,\nonumber \\
(\psi\chi)^\dagger&=&\bar{\chi}_{\dot{\alpha}}\bar{\psi}^{\dot{\alpha}}
=\bar{\chi}\bar{\psi}=\bar{\psi}\bar{\chi}\,. \nonumber
\end{eqnarray}

% [Tutor Note: 原文 PDF 第 23 页，$\sigma^{\mu\nu}$ 定义与旋量标量积。两处易错：(i) 因子 $1/4$（而非某些教材的 $1/2$）来自 $\sigma^{\mu\nu}=\frac14(\sigma^\mu\bar\sigma^\nu-\sigma^\nu\bar\sigma^\mu)$，且第一个式子的指标为 $(\sigma^{\mu\nu})_\alpha^{~\beta}$、第二个为 $(\bar\sigma^{\mu\nu})^{\dot\alpha}_{~\dot\beta}$，缩并时方向不同；(ii) $\psi\chi=\psi^\alpha\chi_\alpha=-\psi_\alpha\chi^\alpha=\chi^\alpha\psi_\alpha=\chi\psi$ 中两个负号分别来自 $\epsilon$ 升指标与 Grassmann 反交换，$({\psi\chi})^\dagger=\bar\chi\bar\psi$ 的次序也不可颠倒。后文所有 Yukawa 项与质量项的符号都依赖这组约定。]

我们再列出一些旋量恒等式
$$
\begin{array}{rl}
\chi\sigma^\mu\bar{\psi} & =-\bar{\psi}\bar{\sigma}^\mu\chi\,, \\
(\chi\sigma^\mu\bar{\psi})^\dagger & =\psi\sigma^\mu\bar{\chi}\,, \\
\chi\sigma^\mu\bar{\sigma}^\nu\psi & =\psi\sigma^\nu\bar{\sigma}^\mu
\chi \,,\\ 
(\chi\sigma^\mu\bar{\sigma}^\nu\psi)^\dagger & =\bar{\psi} 
\bar{\sigma}^\nu \sigma^\mu \bar{\chi}\,.
\end{array}
$$
在上述基底下，狄拉克矩阵以及狄拉克旋量与 Majorana 旋量分别为
\begin{equation}
\gamma^\mu = \left( \begin{array}{cc}
0 & \sigma^\mu \\ \bar{\sigma}^\mu & 0 \end{array} \right)\,, \qquad
\psi_D=\left(\begin{array}{c}\psi_\alpha\\\bar{\chi}^{\dot\alpha}
\end{array}\right)\,,\qquad
\psi_M=\left(\begin{array}{c}\psi_\alpha\\{\bar{\psi}^{\dot{\alpha}}}
\end{array}\right)\,.
\label{sp-comp}
\end{equation}
如通常那样定义 $\gamma_5=-i\gamma^0\gamma^1\gamma^2\gamma^3$。考虑一个沿 $z$ 方向运动的无质量费米子。此时 $P^\mu=
E(1,0,0,1)$，狄拉克方程给出 $(\gamma^0-\gamma^3)\psi=0$。
由于此时螺旋度算符为 $J_3=\frac{i}{2}\gamma^1\gamma^2$，
可得 $J_3\psi=\frac{i}{2}(\gamma^0)^2 \gamma^1\gamma^2\psi =
\frac{i}{2}\gamma^0\gamma^1\gamma^2\gamma^3\psi=-\frac{1}{2}\gamma_5\psi$。
因此，
$$
\gamma_5 =+1 \Rightarrow -ve \;{\rm helicity}\,,
\gamma_5 =-1 \Rightarrow +ve \;{\rm helicity}\,.
$$

% [Tutor Note: 原文 PDF 第 23 页，$\gamma_5$ 与螺旋度。在 (+---) 号差下对 Dirac 方程的解有 $J_3=-\frac12\gamma_5$，故 $\gamma_5=+1$ 对应负螺旋度（原文写作 $+1\Rightarrow -ve$ helicity）。这与 $(-,+,+,+)$ 文献中常见的"$\gamma_5=+1$ 即正螺旋度"相反，引用时务必核对号差约定。]

\subsection{无中心荷的超对称代数}

在没有中心荷时，超对称代数写成
\begin{eqnarray}
&&\{ Q^{I}_{\alpha}, \bar{Q}_{\dot{\alpha} J} \} = 2
\sigma^\mu_{\alpha \dot{\alpha}} P_\mu \delta^{I}_{J}\,, 
\nonumber\\
&&\{Q^I_\alpha , Q^J_\beta\}=0\,,\,\{{\bar Q}_{\dot\alpha I}\,,\, 
{\bar Q}_{\dot\beta J}\}=0\,.
\label{susy-noC}
\end{eqnarray}
这里 $Q$ 与 $\bar Q$ 是超对称生成元，在角动量代数下按自旋 $1/2$ 的算符变换。指标 $I, J$ 从 $1$ 取到 $N$，$N$ 为超对称的总数。此外，超对称生成元与动量算符 $P_\mu$ 对易，从而也与 $P^2$ 对易。因此，代数任一给定表示中的所有态都具有相同的质量。一个理论若要是超对称的，其粒子内容必须构成上述代数的一个表示。超对称代数可以嵌入超 Poincaré 代数，其表示可用 Wigner 方法系统地求得。下面我们简要描述超对称代数的表示。

\noindent\underline{无质量不可约表示}：
\noindent 对无质量态，我们总可转到 $P^{\mu} = M (1,0,0,1)$ 的参照系。此时超对称代数变为
$$
\{ Q^{I}_{\alpha}, \bar{Q}_{\dot{\alpha}J} \} =
\left ( \begin{array}{lr}
0 & 0 \\ 0 & 4M \end{array} \right )
\delta^{I}_{J}\,.
$$
现在，在幺正理论中态的模总是正定的。由于 $Q_\alpha$ 与 $\bar Q_{\dot\alpha}$ 互为共轭，且 $\{ Q_{1}, \bar{Q}_{\dot{1}}\}=0$，由此可知
$Q_1|phys>= \bar{Q}_{\dot{1}}|phys> =0$。至于其余生成元，方便的做法是把它们重新标度为
$$
a^I=\frac{1}{2\sqrt{M}}Q^I_2\,,\qquad
(a^I)^\dagger=\frac{1}{2\sqrt{M}}\bar Q^I_{\dot 2}\,.
$$
于是超对称代数取形式
$$
\{ a^I, (a^J)^\dagger \} = \delta^{IJ}\,,\quad  \{ a^I, a^J \}=0\,,
\quad \{ (a^I)^\dagger, (a^J)^\dagger \} = 0\,.
$$

% [Tutor Note: 原文 PDF 第 23 页（印刷页 22），无质量表示的 $\{Q^I_\alpha,\bar Q_{\dot\alpha J}\}$ 矩阵。取 $P^\mu=M(1,0,0,1)$ 后 $\sigma^\mu P_\mu=\mathrm{diag}(P^0-P^3,\,P^0+P^3)=\mathrm{diag}(0,2M)$，乘以式 (50) 的因子 2 即得右下角的 $4M$。又：由 $\{Q_1,\bar Q_{\dot1}\}=0$ 推出 $Q_1|phys>=0$ 的依据是幺正性，$0=\langle phys|\{Q_1,\bar Q_{\dot1}\}|phys\rangle=\|Q_1|phys\rangle\|^2+\|\bar Q_{\dot1}|phys\rangle\|^2$；注意超荷反对易子为 $+2\sigma^\mu P_\mu$（正号），而 §2.4 超协变导数为 $\{D_\alpha,\bar D_{\dot\alpha}\}=-2i\sigma^\mu\partial_\mu$（负号），两者恰好相反，是本章最常见的符号陷阱。]

这是一个含 $2N$ 个生成元的 Clifford 代数，具有 $2^{N}$ 维表示。从角动量代数的观点看，$a^I$ 是（无质量态）螺旋度的上升算符，$(a^I)^\dagger$ 则是下降算符。我们选取真空，使得 $J_3|\Omega_\lambda>=\lambda 
|\Omega_\lambda>$ 且对一切 $I$ 有 $a^I|\Omega_\lambda>=0$。其余态由 $(a^I)^\dagger$ 作用在真空态上生成。由反对称性可知，含 $m$ 个
$(a^I)^\dagger$、因而螺旋度为 $\lambda -m/2$ 的态，其简并度为 $^N{\bf C}_m$。如此构造的所有态的螺旋度覆盖从 $\lambda$ 到 $\lambda-N/2$ 的范围。例如：
$$
\begin{array}{rrrrrr}
N = 1\,:& \quad |\lambda >,&    |\lambda-1/2>\hphantom{,}& & &\\
N = 2\,:& \quad |\lambda >,& 2\,|\lambda-1/2>,&|\lambda -1>\hphantom{,}&&\\ 
N = 4\,:& \quad |\lambda >,& 4\,|\lambda-1/2>,&
6\,|\lambda -1 >,&4\,|\lambda - 3/2>,& |\lambda -2 >
\end{array}
$$

% [Tutor Note: 原文 PDF 第 23 页。"$a^I$ is a rising operator and $(a^I)^\dagger$ is a lowering operator for the helicity" 一句容易误解：按角动量，$(a^I)^\dagger$ 使 $J_3$ 增加 $1/2$；但原文以"$m$ 个 $(a^I)^\dagger$ 给出螺旋度 $\lambda-m/2$"定义螺旋度，故对螺旋度而言 $(a^I)^\dagger$ 是下降算符。判断升降时请以 $J_3$ 的矩阵元为准，不要只看名称。]

不可约表示不一定具有 CPT 不变性。因此，若要把物理态赋予这些表示，必须补上它们的 CPT 共轭；若某表示是 CPT 自共轭的，则保持不变。下面我们列出补上 CPT 共轭后的表示，并指出可赋予它们的粒子谱：
$$
\begin{array}{lrllllll}
N=1,&\lambda=1/2:&\,\,\, |1/2>\,,&|0>\,,   &|-1/2>\,,&|0>\hphantom{,}&& \\
    &\lambda=1:&\,\,\,|1> \,,&|1/2>\,, &|-1>\,,  &|-1/2>\hphantom{,}&& \\
N=2,&\lambda=1/2:&\,\,\,|1/2>\,,&2|0>\,,&|-1/2>\,,&|-1/2>\,,&2|0>\,,&|1/2>\\
    &\lambda=1:&\,\,\, |1> \,,&2|1/2>\,,&|0>\,,   &|-1>\,, &2|-1/2>\,,&|0>\\
N=4,&\lambda=1:&\,\,\, |1>\,,&4|1/2>\,,&6|0>\,,&4|-1/2>\,,&|-1>\hphantom{,}&
\end{array}
$$

% [Tutor Note: 原文 PDF 第 24 页（印刷页 23）。"we have to suplement them with their CPT conjugates" 与 "which can be assiged to them" 中 "suplement"、"assiged" 均为原文拼写错误，应作 "supplement"、"assigned"。又：表中 $N=2,\ \lambda=1/2$ 一行把 $|1/2>,2|0>,|-1/2>$ 与其 CPT 共轭 $|-1/2>,2|0>,|1/2>$ 并列，故 $|0>$、$|1/2>$ 各出现两次，这不是重复排版，而是把表示扩充为 CPT 自共轭的结果。]

于是，对 $N=1$，若 $\lambda=1/2$，表示包含一个 Majorana 旋量与一个复标量（标量多重态）；若 $\lambda=1$，则包含一个无质量矢量与一个 Majorana 旋量（矢量多重态）。对 $N=2$ 且 $\lambda=1/2$，有两个 Majorana 旋量（或一个狄拉克旋量）和两个复标量；该表示的粒子内容与两份 $N=1$、$\lambda=1/2$ 多重态相同。对 $N=2$ 且 $\lambda=1$，有一个无质量矢量、两个 Majorana 旋量和一个复标量。注意该多重态的粒子内容恰好等于 $\lambda=1/2$ 与 $\lambda=1$ 的两个 $N=1$ 多重态之和。对 $N=4$，表示是自共轭的，容纳一个无质量矢量、两个狄拉克费米子和三个复标量。

\noindent\underline{有质量不可约表示}：
\noindent 对有质量态，我们总可转到静止系 $P_\mu=(M,0,0,0)$，并定义
$$
a^I_\alpha=Q^I_\alpha/\sqrt{2M}\,,\quad (a^I_\alpha)^\dagger= 
{\bar Q}_{\dot\alpha I}/\sqrt{2M}\,.
$$
于是超对称代数化为
$$
\{a^I_1,(a_1^J)^\dagger\}=\delta^{IJ}\,,\qquad 
\{a^I_2,(a_2^J)^\dagger\}=\delta^{IJ}\,,
$$
其余反对易子全为零。Clifford 真空由 $a^I_\alpha |\Omega>=0$ 定义，表示通过把 $(a^I_\alpha)^\dagger$ 作用到 $\Omega$ 上构造。设
$|\Omega>$ 为自旋单态，则第 $m$ 层有 $^{2N}{\rm C}_m$ 个态，表示的维数为
$\sum_{m=0}^{2N}{}^{2N}{\rm\bf C}_m=2^{2N}$。所能达到的最大自旋是 $N/2$，而不是人们朴素预期的 $N$。原因是
$(a^I_1)^\dagger(a^I_2)^\dagger=\frac{1}{2}
\epsilon^{\alpha\beta} (a^I_\alpha)^\dagger(a^I_\beta)^\dagger$ 是标量。因此 $m=2N$ 的态与真空一样自旋为零。具有给定自旋的态的简并度由群 $USp(2N)$ 的不可约表示标记，这里不作讨论。作为替代，我们看最简单的例子。对 $N=1$，有质量表示包含 $2^2=4$ 个态，
$$
|\Omega>,\quad a^\dagger_\alpha |\Omega>,\quad\frac{1}{\sqrt{2}}
\epsilon^{\alpha \beta} a^\dagger_\alpha a^\dagger_\beta |\Omega>\,,
$$
其自旋内容为 $(0)\oplus (1/2)\oplus (0)$。这里 $(j)$ 表示总自旋为 $j$、简并度为 $2j+1$ 的态。因此在上例中，我们有一个外尔（或 Majorana）旋量和一个复标量
$(\lambda , \phi)$。对 $N=2$，表示包含 $2^4=16$ 个态，它们在角动量的 $SU(2)$ 下分解为
$5(0)\oplus 4(1/2)\oplus 1(1)$。$N=4$ 的有质量多重态有
$2^8=256$ 个态，并包含一个自旋 $2$ 的态。

% [Tutor Note: 原文 PDF 第 25 页（印刷页 24）。三处提醒：(i) 最大自旋为 $N/2$ 而非 $N$，关键是 $(a^I_1)^\dagger(a^I_2)^\dagger$ 在自旋指标上反对称、只能给出标量（原文用 $\frac12\epsilon^{\alpha\beta}(a^I_\alpha)^\dagger(a^I_\beta)^\dagger$ 说明）；(ii) $USp(2N)$ 为紧致辛群，有质量多重态的简并度按其不可约表示分类，原文明确略去不证；(iii) "the dimension of the the representation" 中衍出一个 "the"，系原文笔误。]

至此我们考虑的都是以单态真空为基础的表示。现在考虑自旋为 $j$、简并度为 $2j+1$ 的真空 $|\Omega_j>$。此时表示含 $(2j+1)2^{2N}$ 个态。其谱可通过把 Clifford 代数的 $j=0$ 表示与自旋 $j$ 按角动量相加规则组合而得到。例如，要得到以
$|\Omega_j>$ 为基础的 $N=1$ 表示，我们把 $(0)\oplus(1/2)\oplus (0)$ 与 $(j)$ 组合，得到 $(j)\oplus(j+1/2)\oplus(j-1/2)\oplus(j)$。对 $j=1/2$，得到
$(1/2)\oplus(1)\oplus(0)\oplus(1/2)$，对应于一个规范场、一个狄拉克费米子和一个标量场，三者质量相同。注意在所有这些情形中，玻色与费米的自由度数都相同。

\subsection{含中心荷的超对称代数}

如 Haag、Lapuszanski 与 Sohnius \cite{HLS} 所示，超对称代数 (\ref{susy-noC}) 容许中心扩张，可推广为
\begin{eqnarray}
&&\{ Q^{I}_{\alpha}, \bar{Q}_{\dot{\beta}J} \} =
2 \sigma^\mu_{\alpha \dot{\beta}} P_\mu \delta^{I}_{J}\,, 
\nonumber\\
&&\{ Q^I_\alpha, Q^J_\beta \}=2{\sqrt 2}\epsilon_{\alpha\beta} Z^{IJ}\,,
\nonumber \\
&&\{\bar{Q}_{\dot{\alpha}I},\bar{Q}_{\dot{\beta}J} \} =
2{\sqrt 2}\epsilon_{\dot{\alpha} \dot{\beta}} Z^{*}_{IJ}\,,
\label{susy-C}
\end{eqnarray}
其中 $Z$ 与 $Z^*$ 为在 $I$ 与 $J$ 上反对称的中心荷矩阵。我们先讨论 $N$ 为偶数的情形。借助一个幺正变换，可把 $Z$ 化为斜对角形式：$\tilde{Z}^{IJ} = U^{I}_{A}U^{J}_{B}Z^{AB}$，使其形如 $Z=\epsilon\otimes D$，其中 $D$ 是 $N/2$ 维对角矩阵。于是，标记超对称个数的指标 $I$ 可分解为 $(a,m)$：$a=1,2$ 来自反对称张量 $\epsilon$，$m=1,...,N/2$ 来自对角矩阵 $D$。再作一次手征转动，可使 $D$ 的本征值取实数。一旦完成斜对角化，只需考虑 $N=2$ 超对称就够了；此时代数取形式
\begin{eqnarray}
&&\{Q^a_\alpha,{\bar Q}_{\dot\beta b}\}=2(\sigma^\mu)_{\alpha\dot\beta}  
P_\mu \delta^{a}_{b}\,,\nonumber \\
&&\{ Q^a_\alpha,Q^b_\beta\}=2{\sqrt 2}\epsilon_{\alpha\beta}
\epsilon^{ab}Z\,, \nonumber \\ 
&&\{{\bar Q}_{\dot\alpha a},{\bar Q}_{\dot\beta b}\}= 
2{\sqrt 2}\epsilon_{\dot\alpha\dot\beta}\epsilon_{ab} Z \,.
\label{Susy-C'}
\end{eqnarray}

% [Tutor Note: 原文 PDF 第 26 页（印刷页 25），式 (51) 前后。两处易错与一处笔误：(i) 中心荷只出现在 $\{Q,Q\}$ 与 $\{\bar Q,\bar Q\}$ 中，系数为 $2\sqrt2\,\epsilon_{\alpha\beta}Z^{IJ}$（不是 $2\epsilon Z$），$\{Q,\bar Q\}$ 中没有中心荷；(ii) 斜对角化后 $Z=\epsilon\otimes D$，指标 $(a,m)$ 的拆分意味着不同 $m$ 的子代数相互对易，"只需考虑 $N=2$"的根据在此——每个 $m$ 给出独立的 $N=2$ 代数，各自给出 $M\geq\sqrt2|Z_m|$，故合起来要求 $M\geq\sqrt2\max_m|Z_m|$；(iii) 原文 "Haag, Lapuszanski and Sohnius" 中 "Lapuszanski" 应作 "Łopuszański"（即 Haag--Łopuszański--Sohnius 定理，文献 \cite{HLS}）。]

由于 $Z$ 与所有生成元对易，我们可以把它固定为该给定表示的本征值。现在定义：
$$
a_\alpha=\frac{1}{2} \{ Q^1_\alpha +
\epsilon_{\alpha\beta} (Q^2_\beta)^\dagger \}\,, \qquad
b_\alpha=\frac{1}{2} \{ Q_\alpha^1 -
\epsilon_{\alpha\beta} (Q^2_\beta)^\dagger \}\,. 
$$
于是代数 (\ref{susy-C}) 化为
\begin{equation}
\{a_\alpha,a^\dagger_\beta\}=\delta_{\alpha\beta}(M+{\sqrt 2}Z)\,,\qquad
\{ b_\alpha,b^\dagger_\beta\}=\delta_{\alpha\beta} (M-{\sqrt 2}Z)\,,
\label{susy-ab}
\end{equation}
其余反对易子全为零。

% [Tutor Note: 原文 PDF 第 26 页（印刷页 25），式 (53) 及其前一句。两点：(i) 原文写 "Then, the algebra (51) reduces to"，但按上下文此处应指斜对角化后得到的 $N=2$ 代数式 (52)（式 (51) 是带 $I,J$ 指标的完整代数），疑为原文式号笔误，译文照原文译出，请教授核对；(ii) 组合 $a_\alpha=\frac12\{Q^1_\alpha+\epsilon_{\alpha\beta}(Q^2_\beta)^\dagger\}$ 中的 $(Q^2_\beta)^\dagger$ 是 $\bar Q_{\dot\beta}^{\,2}$，其指标结构（带点指标 $\dot\beta$ 与 $a_\alpha$ 的 $\alpha$ 通过 $\epsilon$ 相配）是得到 $\delta_{\alpha\beta}(M\pm\sqrt2Z)$ 的关键，符号写反则长短多重态的判据会颠倒。]

由于所有物理态都有正定模，可知对无质量态，中心荷只能平庸地实现（即 $Z=0$）。对有质量态，这给出质量界 $M\geq {\sqrt 2}|Z|$。当
$M={\sqrt 2}|Z|$ 时，(\ref{susy-ab}) 中有一组算符平庸实现，代数与无质量情形相似，表示的维数大大减小。例如，约化的有质量 $N = 2$ 多重态与无质量 $N = 2$ 多重态具有相同数目的态。因此，带中心荷的 $N=2$ 代数的表示可分为长多重态（$M >
{\sqrt 2}|Z|$）与短多重态（$M={\sqrt 2}|Z|$）。

质量界 $M\geq {\sqrt 2}|Z|$ 使人想起 Georgi–Glashow 模型中的 Bogomol'nyi 界。事实上，在 Georgi–Glashow 模型的超对称版本（基于无中心荷的代数）中，孤子解确实在超对称代数中产生中心扩张项，从而实现 (\ref{susy-C})\cite{Olive-Witten}。中心荷的起源很容易理解：超对称荷 $Q$ 与 $\bar Q$ 是场中局域表达式（超流的时间分量）的空间积分；计算它们的反对易子时，会遇到通常被忽略的表面项。然而，在存在电荷与磁荷时，这些表面项不为零，从而给出中心荷。正如本节末尾将要显式证明的，结果为
\begin{equation}
Z = a (q+ig)= ae (n+m\tau)\,,
\label{Ztau}
\end{equation}
于是 $M\geq {\sqrt 2}|Z|$ 与 Bogomol'nyi 界 (\ref{Bbound}) 一致。由 (\ref{susy-ab}) 可以看出，BPS 态（即达到该界的态）被一半超对称生成元湮灭，因而属于
(\ref{susy-C}) 的约化表示。这一点的重要推论是：对 BPS 态，其荷与质量之间的关系由超对称决定，在量子理论中不受微扰或非微扰修正。原因在于，若该关系被修改，这些态就不再属于短多重态。另一方面，量子修正不应产生把短多重态变成长多重态所需的额外自由度。既然没有其他可能，我们便得出结论：对短多重态，关系 $M={\sqrt 2}|Z|$ 既不受到微扰修正，也不受到非微扰修正。

% [Tutor Note: 原文 PDF 第 27 页（印刷页 26），式 (54)。这里 $a=|\langle\phi\rangle|$ 是 §1 中希格斯真空的模（原文 PDF 第 21 页明确取 $|\phi|=a$），$e$ 为基本电荷，$n,m$ 分别为电荷与磁荷的整数量子数，$\tau=\theta/2\pi+4\pi i/e^2$（原文 PDF 第 22 页式 (47)）；式 (54) 把中心荷归一化为 $ae(n+m\tau)$，故 $M\geq\sqrt2|Z|$ 与 Bogomol'nyi 界 (35) 同义。"量子修正不应产生额外自由度"一句是原文的启发式论证（短多重态的自由度数由代数决定，不能连续改变），并非严格证明。]

\subsection{$N=1$ 超对称的局域表示}

本小节描述超对称在量子场论中如何作用于局域场。众所周知，Poincaré 群自然地作用在时空坐标上；其他一切对象都按定义在时空流形上的张量或旋量的分量变换。类似地，超对称变换自然地作用在时空的一个扩充上，这个扩充称为“超空间”。于是量子场按定义在超空间上的“超场”的分量变换。下面我们先描述这些概念，再引入手征超场与矢量超场。

\noindent\underline{超空间}：
\noindent 超空间是在时空坐标 $x^\mu$ 上添加四个旋量自由度 $\theta^\alpha,\bar\theta_{\dot\alpha}$ 得到的。旋量指标用 $\epsilon$ 张量升降，且 $\theta\theta = \theta^\alpha\theta_\alpha = -2
\theta^1\theta^2$。类似地，$\bar\theta\bar\theta=
\bar\theta_{\dot\alpha}\bar\theta^{\dot\alpha}= 2\bar\theta_{\dot 1} 
\bar\theta_{\dot 2}$。我们还有
$$
\theta^{\alpha}\theta^{\beta}=\frac{1}{2} \epsilon^{\alpha \beta}
\theta \theta\,,\qquad
\bar{\theta}_{\dot{\alpha}}\bar{\theta}_{\dot{\beta}}=-\frac{1}{2}
\epsilon_{\dot{\alpha}\dot{\beta}} \bar{\theta}\bar{\theta}\,,\qquad
\theta \sigma^\mu \bar{\theta} \theta \sigma^\nu \bar{\theta}=
\frac{1}{2}\theta \theta \bar{\theta}\bar{\theta} \eta^{\mu\nu}\,.
$$
这些公式是 Fierz 重排的基础。

% [Tutor Note: 原文 PDF 第 27 页（印刷页 26），超空间约定。注意 $\theta\theta\equiv\theta^\alpha\theta_\alpha=-2\theta^1\theta^2$ 与 $\bar\theta\bar\theta\equiv\bar\theta_{\dot\alpha}\bar\theta^{\dot\alpha}=+2\bar\theta_{\dot1}\bar\theta_{\dot2}$ 的符号相反：这源于 $\theta$ 与 $\bar\theta$ 指标升降所用的同一 $\epsilon$ 张量的缩并方式；换用其他教材的 $\epsilon$ 约定时，凡 $\theta\theta$ 型的项都会整体变号。第三式 $\theta\sigma^\mu\bar\theta\,\theta\sigma^\nu\bar\theta=\frac12\theta\theta\bar\theta\bar\theta\,\eta^{\mu\nu}$ 是 Fierz 重排的基本恒等式，其系数 $1/2$ 与 $\eta^{\mu\nu}$ 的位置都不可改动。]

在 $N=1$ 的超对称变换 (\ref{susy-noC}) 下（变换参数记为 $\xi$ 与 $\bar\xi$），超空间坐标取如下变换：
\begin{eqnarray}
x^\mu &\rightarrow& x'^\mu=x^\mu + i\theta\sigma^\mu\bar\xi
- i\xi\sigma^\mu\bar\theta \,,\nonumber\\
\theta &\rightarrow&\theta'=\theta + \xi\,, \nonumber\\
\bar\theta &\rightarrow& \bar\theta'=\bar\theta +\bar\xi\,.
\label{superspace}
\end{eqnarray}
由于这些变换由算符
$\xi^\alpha Q_\alpha + \bar{\xi}_{\dot\alpha}\bar{Q}^{\dot\alpha}$ 实现，我们可以容易地得到超荷作用在超空间上的表示
\begin{equation}
Q_\alpha=\frac{\partial}{\partial\theta^\alpha}- 
i\sigma^\mu_{\alpha\dot{\alpha}}\bar{\theta}^{\dot{\alpha}}\,
\partial_\mu\,,\qquad \bar{Q}_{\dot\alpha} = -\frac{\partial}{\partial
\bar{\theta}^{\dot\alpha}}+i\theta^\alpha 
\sigma^\mu_{\alpha\dot{\alpha}}\, \partial_\mu\,.
\label{superspace-Q}
\end{equation}
它们满足 $\{Q_\alpha ,\bar{Q}_{\dot\alpha}\}= 
2i\sigma^\mu_{\alpha\dot\alpha}\,\partial_\mu$。此外，利用链式法则容易看出，$\partial/\partial x^\mu$ 在 (\ref{superspace}) 下不变，而 $\partial/\partial\theta$ 与 $\partial/\partial\bar\theta$ 则不然。因此我们引入超协变导数
\begin{equation}
D_\alpha = \frac{\partial}{\partial \theta^\alpha} +
i\sigma^\mu_{\alpha\dot\alpha}\,\partial_\mu \,,\qquad
\bar{D}_{\dot{\alpha}} = -\frac{\partial}{\partial
\bar{\theta}^{\dot{\alpha}}}-i\sigma^\mu_{\alpha \dot{\alpha}}
\theta^{\alpha} \, \partial_\mu \,.
\label{superspaceD}
\end{equation}
它们满足 $\{ D_\alpha,\bar{D}_{\dot\alpha}\}=-2i\sigma^\mu_{\alpha
\dot{\alpha}}\,\partial_\mu$，并与 $Q$、$\bar Q$ 对易。

% [Tutor Note: 原文 PDF 第 28 页（印刷页 27），式 (57) 有一处原文公式笔误：$D_\alpha=\partial/\partial\theta^\alpha+i\sigma^\mu_{\alpha\dot\alpha}\partial_\mu$ 中漏写了 $\bar\theta^{\dot\alpha}$ 因子，正确形式应为 $D_\alpha=\partial/\partial\theta^\alpha+i\sigma^\mu_{\alpha\dot\alpha}\bar\theta^{\dot\alpha}\partial_\mu$；否则 $\{D_\alpha,\bar D_{\dot\alpha}\}$ 只能得到 $-i\sigma^\mu_{\alpha\dot\alpha}\partial_\mu$ 项（只有正确形式才给出原文所写的 $-2i\sigma^\mu\partial_\mu$）。正文公式按原文照录，正确形式在此注明。另需注意 $\{Q,\bar Q\}=+2i\sigma^\mu\partial_\mu$ 与 $\{D,\bar D\}=-2i\sigma^\mu\partial_\mu$ 差一负号，且 $Q_\alpha$ 与 $\bar Q_{\dot\alpha}$ 的一阶项符号相反（分别对应式 (55) 中 $x^\mu$ 位移的两项），这三处合起来是本章最易出错的地方。]

\noindent\underline{超场}：
\noindent 超场是超空间上的函数，记作
$F(x,\theta,\bar\theta)$。由于 $\theta$ 坐标反交换，最一般的 $N=1$ 超场总可展开为
\begin{eqnarray}
F(x,\theta,\bar\theta)&=& f(x)+\theta\phi(x) + \bar\theta\bar\chi(x)
+\theta\theta m(x) +\bar\theta\bar\theta n(x) + 
\theta\sigma^\mu\bar\theta v_\mu (x) \nonumber\\
&+&\theta\theta\bar\theta\bar\lambda(x)+\bar\theta\bar\theta\theta\psi(x)
+ \theta\theta\bar\theta\bar\theta d(x)\,. 
\label{superfield}
\end{eqnarray}
显然，超场的任意函数本身仍是超场。在超对称下，超场按 $\delta F=(\xi
Q+\bar\xi\bar Q)F$ 变换，由此可得各分量场的变换。注意，由于 $d(x)$ 是多重态中量纲最高的分量，它在超对称下的变分总是其他分量的全导数。因此，忽略表面项，该分量的时空积分在超对称下不变。这告诉我们：超对称拉格朗日密度可以构造为某个适当超场的最高量纲分量。为描述物理系统，我们并不需要超场的全部分量；对超场施加适当约束即可挑出相关的分量。

% [Tutor Note: 原文 PDF 第 28 页，式 (58) 的超场展开。系数 $1$（$f$）、$1$（$\theta\phi$）、$1$（$\theta^2m$）、$1/2$ 量的 $\theta^2\bar\theta^2d$ 都是相对于 $\theta\theta=\theta^\alpha\theta_\alpha$ 的约定；注意展开中 $\theta^2$ 项、$\bar\theta^2$ 项彼此独立，$\theta\sigma^\mu\bar\theta v_\mu$ 是唯一同时含 $\theta$ 与 $\bar\theta$ 的双线性项。$d(x)$ 的变分是全导数这一事实，是"超对称拉氏量 = 超场最高分量"这一构造原则的根据。]

\noindent\underline{手征多重态}：
\noindent $N=1$ 标量多重态由一个带约束的超场表示：
$$
\bar{D}_{\dot{\alpha}} \Phi = 0 \,.
$$
这样的超场称为手征超场。注意，若取
$y^\mu = x^\mu +i\theta\sigma^\mu\bar{\theta}$，则
$$
\bar{D}_{\dot\alpha}y^\mu=0,\quad\bar{D}_{\dot\alpha}\theta^\beta=0\,.
$$
因此 $(y,\theta)$ 的任意函数都是手征超场。可以证明这也是必要条件。于是任意手征超场可展开为
\begin{equation}
\Phi(y,\theta)=A(y)+\sqrt{2}\theta\psi(y)+\theta\theta F(y)\,.
\label{csfy}
\end{equation}
这里 $A$ 与 $\psi$ 分别是标量分量与费米分量，$F$ 是为使代数在壳外闭合所需的辅助场。类似地，反手征超场由 $D_{\alpha}\Phi^\dagger=0$ 定义，可展开为
\begin{equation}
\Phi^\dagger(y^\dagger,\bar\theta) =A^\dagger(y^\dagger)+\sqrt{2}
\bar{\theta}\bar{\psi}(y^\dagger)+\bar{\theta}\bar{\theta} 
F^\dagger(y^\dagger)\,,
\label{acsfy}
\end{equation}
其中 $y^{\mu\dagger}=x^\mu-i\theta\sigma^\mu\bar{\theta}$。
手征超场之积仍是手征超场。一般地，手征超场的任意函数都是手征超场：
\begin{eqnarray}
{\cal W}(\Phi_{i}) &=& {\cal W}(A_{i} + \sqrt{2} \theta \psi_{i} +
\theta \theta F_{i} ) \nonumber \\
&=&{\cal W}(A_{i}) + \frac{\partial {\cal W}}{\partial A_{i}}
\sqrt{2} \theta \psi_{i} + \theta \theta \left (\frac{\partial 
{\cal W}}{\partial A_{i}} F_{i} - \frac{1}{2} \, \frac{\partial^{2}
{\cal W}}{\partial A_i A_j}\, \psi_{i} \psi_{j} \right)\,.
\label{spotential}
\end{eqnarray}
${\cal W}$ 称为超势。用原始变量表示，$\Phi$ 与 $\Phi^\dagger$ 取形式
\begin{eqnarray}
\Phi(x,\theta,\bar\theta) &=& A(x)+i\theta\sigma^\mu \bar{\theta}
\partial_\mu A-\frac{1}{4}\theta^2 \bar{\theta}^2 \Box A 
\nonumber\\
&&+\sqrt{2}\theta\psi(x)-\frac{i}{\sqrt{2}}\theta\theta
\partial_\mu\psi\sigma^\mu\bar{\theta} +\theta \theta F(x)\,,
\label{csfield}
\\
\Phi^\dagger(x,\theta,\bar\theta)&=&A^\dagger(x)-i\theta\sigma^\mu
\bar{\theta} \partial_\mu A^\dagger - \frac{1}{4} \theta^{2}
\bar{\theta}^{2} \Box A^\dagger \nonumber\\
&&+\sqrt{2} \bar{\theta} \bar{\psi} (x) + \frac{i}{\sqrt{2}} \,
\bar{\theta} \bar{\theta} \, \theta \sigma^\mu \, \partial_\mu
\bar{\psi} + \bar{\theta} \bar{\theta} F^\dagger (x)\,.
\label{asfield}
\end{eqnarray}

% [Tutor Note: 原文 PDF 第 28 页（印刷页 27），式 (59) 及其下方。原文 "Here, $A$ and $\psi$ are the fermionic and scalar components respectively" 中 "fermionic" 与 "scalar" 的词序写反了：由式 (59) 可见 $A$ 是标量分量、$\psi$ 是费米分量。译文按正确物理含义译出，原文错误在此注明。另外，式 (59) 中 $A,\psi,F$ 都依赖 $y=x+i\theta\sigma\bar\theta$ 而非 $x$，这正是手征超场的本质（$\bar D\Phi=0$ 的解全部集中在 $(y,\theta)$ 上）；式 (62)--(63) 换回 $(x,\theta,\bar\theta)$ 后出现的 $\Box A$、$\partial_\mu\psi$ 等项正是 $y$ 与 $x$ 之差的高阶效应，作分量计算时不可遗漏。]

\noindent\underline{矢量多重态}：
\noindent 该多重态由满足 $V=V^\dagger$ 的实超场表示。按分量写，其形式为
$$
\begin{array}{lll}
V(x, \theta, \bar{\theta}) =& C + i \theta  \chi - i \bar{\theta}
\bar{\chi} + \frac{i}{2} \theta^{2} (M + i N) - \frac{i}{2}
\bar{\theta}^{2} (M - iN) \\
& - \theta \sigma^\mu \bar{\theta} A_\mu + i \theta^{2} \bar{\theta}
(\bar{\lambda} + \frac{i}{2} \bar{\sigma}^\mu \partial_\mu \chi )\\
& - i \bar{\theta}^{2} \theta (\lambda + \frac{i}{2} \sigma^\mu
\partial_\mu \bar{\chi} ) + \frac{1}{2} \theta^{2} \bar{\theta}^{2}
(D - \frac{1}{2} \Box C)\,.
\end{array}
$$
其中许多分量可用阿贝尔规范变换 $V \rightarrow V + \Lambda + \Lambda^\dagger$ 规范掉，这里
$\Lambda$（$\Lambda^\dagger$）是手征（反手征）超场。在所谓 Wess–Zumino 规范中取 $C=M=N=\chi= 0$，于是
$$
V=-\theta \sigma^\mu\bar{\theta}A_\mu+i\theta^{2}
\bar{\theta}\bar{\lambda} - i \bar{\theta}^{2} \theta \lambda +
\frac{1}{2} \theta^{2} \bar{\theta}^{2} D\,.
$$
在此规范下，$V^2=\frac{1}{2}A_\mu A^\mu \theta^{2}\bar{\theta}^{2}$，且 $V^{3}=0$。Wess–Zumino 规范破坏超对称，但不破坏阿贝尔规范场 $A_\mu$ 的规范对称性。阿贝尔场强定义为
$$
W_{\alpha} = - \frac{1}{4} \bar{D}^{2} D_{\alpha} V\,,\qquad
\bar{W}_{\dot{\alpha}}=-\frac{1}{4}D^{2} \bar{D}_{\dot{\alpha}}V\,.
$$
$W_{\alpha}$ 是手征超场。由于它是规范不变的，可在 Wess–Zumino 规范中计算，其形式为
\begin{equation}
W_\alpha =-i\lambda_\alpha (y)+\theta_\alpha D-\frac{i}{2} 
(\sigma^\mu\bar{\sigma}^\nu \theta)_\alpha \, F_{\mu\nu}+ 
\theta^2 (\sigma^\mu \partial_\mu \bar{\lambda} )_{\alpha}\,,
\label{W-abelian}
\end{equation}
其中 $F_{\mu\nu}=\partial_\mu A_\nu -\partial_\nu A_\mu$ 是熟知的阿贝尔场强张量。

% [Tutor Note: 原文 PDF 第 29 页（印刷页 28），矢量超场与 Wess--Zumino 规范。四点是核对重点：(i) $V$ 的展开系数 $i/2$（$\theta^2$ 项）、$-i/2$（$\bar\theta^2$ 项）、负号（$\theta\sigma^\mu\bar\theta A_\mu$ 项）由 $V=V^\dagger$ 的实性条件固定，改动任一系数都会破坏实性；(ii) 规范变换为 $V\to V+\Lambda+\Lambda^\dagger$ 且 $\Lambda$ 必须是手征超场，这是保持 $V$ 为实超场且 $W_\alpha$ 手征的前提；(iii) 只有在这一规范下才有 $V^3=0$ 与 $V^2=\frac12A_\mu A^\mu\theta^2\bar\theta^2$，超场多项式才能截断；(iv) 规范固定条件 $C=M=N=\chi=0$ 与超对称不相容（WZ 规范破坏超对称但不破坏规范对称性），故按此规范作超对称变换时必须配合补偿变换；原文明确说明这一点，切勿把 WZ 规范下的分量变换直接当作超对称变换。]

在非阿贝尔情形，$V$ 属于规范群的伴随表示：$V=V_AT^A$，其中 $T^{A\dagger}=T^A$。规范变换现在由
$$
e^{-2V} \rightarrow e^{-i \Lambda^\dagger} e^{-2V} e^{i \Lambda}
\qquad {\rm where}, \; \Lambda = \Lambda_{A}T^{A}
$$
实现。非阿贝尔规范场强定义为
$$
W_{\alpha} = \frac{1}{8} \bar{D}^{2} e^{2V} D_{\alpha} e^{-2V}
$$
并按
$$
W_\alpha\rightarrow W'_\alpha= e^{-i\Lambda}W_\alpha e^{i\Lambda}\,.
$$
变换。按分量写，其形式为
\begin{equation}
W_\alpha =T^a \left (-i \lambda^{a}_{\alpha} + \theta_{\alpha}
D^{a} - \frac{i}{2}(\sigma^\mu \bar{\sigma}^\nu \theta )_{\alpha}
F^{a}_{\mu\nu}+\theta^2 \sigma^\mu D_\mu\bar{\lambda}^a \right)\,, 
\label{W-nonabelian}
\end{equation}
其中
$$
F^a_{\mu\nu}=\partial_\mu A^a_\nu -\partial_\nu A^a_\mu + 
\, f^{abc} A^b_\mu A^c_\nu\,,\qquad
D_\mu\bar{\lambda}^a=\partial_\mu\bar{\lambda}^a +
\,f^{abc} A_\mu^b \bar{\lambda}^c \,.
$$
下一节我们将用超场构造超对称拉格朗日量。

% [Tutor Note: 原文 PDF 第 29--30 页（印刷页 28--29），场强的超场定义。注意阿贝尔情形 $W_\alpha=-\frac14\bar D^2D_\alpha V$（因子 $1/4$），非阿贝尔情形 $W_\alpha=\frac18\bar D^2 e^{2V}D_\alpha e^{-2V}$（因子 $1/8$）：两者归一化不同，比较系数时务必分清。非阿贝尔规范变换为 $e^{-2V}\to e^{-i\Lambda^\dagger}e^{-2V}e^{i\Lambda}$（左乘 $\Lambda^\dagger$、右乘 $\Lambda$），与常见的 $e^{2V}$ 写法的符号约定相反；式 (65) 中 $f^{abc}A^b_\mu A^c_\nu$ 与 $D_\mu\bar\lambda^a=(\partial_\mu+f^{abc}A^b_\mu)\bar\lambda^c$ 的符号由 $T^a$ 的伴随表示约定固定，换用 $T^a_{bc}=+if^{abc}$ 的教材时整体反号，需逐项核对。]

\subsection{$N=1$ 拉格朗日量的构造}

本节将构造标量多重态与矢量多重态的 $N=1$ 拉格朗日量。它们是我们真正关心的 $N=2$ 拉格朗日量的基本构件。如前所述，超对称拉格朗日量可构造为超场的最高分量；于是问题化为寻找合适的超场。

\noindent\underline{标量多重态的拉格朗日量}：
\noindent 先考虑手征超场与反手征超场之积
$\Phi_i^\dagger\Phi_j$。这是一般超场，其最高分量可用
(\ref{csfield}) 与 (\ref{asfield}) 算出为
\begin{eqnarray}
\Phi^\dagger_i\Phi_j\mid_{\theta^{2}\bar{\theta}^{2}}\,\,\,=
&-&\frac{1}{4}\,A^\dagger_i\Box A_j -\frac{1}{4}\Box A^\dagger_i 
A_j+F^\dagger_i F_j +\frac{1}{2}\,\partial_\mu A^\dagger_i\partial^\mu
A_j \nonumber\\
&-& \frac{i}{2}\,\psi_j\sigma^\mu \partial_\mu \bar{\psi}_{i} +
\frac{i}{2} \partial_\mu \psi_j \sigma^\mu\bar{\psi}_i\,.\nonumber
\end{eqnarray}
略去若干全导数并对 $i=j$ 求和，得到自由场拉格朗日量
$$
{\cal L}=
\Phi^\dagger_i \Phi_i \mid_{\theta^2\bar\theta^2}
=\partial_\mu A^\dagger_i \partial^\mu A_i + F^\dagger_i F_i -i
\bar{\psi}_i \bar{\sigma}^\mu \partial_\mu \psi_i\,.
$$
这是一个无质量标量与一个无质量费米子的自由拉格朗日量，其中的辅助场可用其运动方程消去。超对称相互作用项可用超势 (\ref{spotential}) 及其共轭构造，二者分别是 $\Phi$ 与 $\Phi^\dagger$ 的全纯函数。此外注意，场 $\Phi$ 的空间可带有非平凡度规 $g^{ij}$；此时标量动能项例如取形式 $g^{ij}\partial_\mu
A_i^\dagger \partial^\mu A_j$，其他项也作相应修改。在此情形下，上述自由场拉格朗日量必须换成一个非线性
$\sigma$ 模型。于是，标量多重态（含相互作用项）的最一般的 $N=1$ 超对称拉格朗日量为
$$
{\cal L}=\int\,d^{4}\theta\,K(\Phi,\Phi^\dagger) +
\int d^{2} \theta {\cal W}(\Phi) + \int d^{2} \bar{\theta} 
\bar{\cal W}(\Phi^\dagger)\,.
$$
注意 $\theta$ 积分取出超场的最高分量；在我们的约定下，$\int d^2 \theta \theta^2=1$ 且
$\int d^2 \bar\theta \bar\theta^2=1$。用非全纯函数 $K(A, A^\dagger)$ 表示，场空间上的度规为
$g^{ij}=\partial^2K/\partial A_i\partial A^\dagger_j$。正因如此，函数 $K(\Phi,\Phi^\dagger)$ 被称为凯勒势。

% [Tutor Note: 原文 PDF 第 30 页（印刷页 29），$\Phi_i^\dagger\Phi_j$ 的最高分量与自由拉格朗日量。核对要点：(i) $-\frac14A_i^\dagger\Box A_j-\frac14\Box A_i^\dagger A_j$ 两项经分部积分合成 $+\frac12\partial_\mu A_i^\dagger\partial^\mu A_j$（实标量动能正确）；(ii) 费米项 $-\frac i2\psi_j\sigma^\mu\partial_\mu\bar\psi_i+\frac i2\partial_\mu\psi_j\sigma^\mu\bar\psi_i$ 分部积分后化为 $-i\bar\psi_i\bar\sigma^\mu\partial_\mu\psi_i$，相对符号极易写反，可用厄米性校验；(iii) $F_i^\dagger F_j$ 为辅助场项，原文 "summing over $i=j$" 系笔误，意思应为"取 $i=j$ 并对 $i$ 求和"。又：原文此处把 $g^{ij}$ 记作场空间度规（通常教材写作 $g_{i\bar\jmath}$），本译本按原记号照录。]

对可重整理论，$K$ 与 ${\cal W}$ 的形式并非任意，而要受 $R$ 对称性约束。该对称性作用在手征超场上如下
\begin{eqnarray}
R\,\Phi (x,\theta) = \Phi'(x,\theta) &=& e^{2 in\alpha}
\Phi (x, e^{-i\alpha}\theta)\,, \nonumber \\
R\,\Phi^\dagger(x,\bar\theta ) = \Phi'^\dagger(x,\theta) &=& e^{-2
in\alpha}\Phi^\dagger(x, e^{i \alpha}\bar{\theta})\,. \nonumber
\end{eqnarray}
在此之下，各分量场按如下方式变换：
$$
\begin{array}{lll}
A & \rightarrow e^{2i n \alpha} A \,,\\
\psi & \rightarrow e^{2i(n-1/2) \alpha} \psi\,, \\
F & \rightarrow e^{2i(n-1) \alpha} F\,.
\end{array}
$$
我们称 $n$ 为 $R$ 特征标。由于 $\theta \rightarrow
e^{+i\alpha}\theta$，即 $d^2\theta \rightarrow e^{-2 i\alpha}d^2
\theta$，${\cal W}$ 每一项中各超场的 $R$ 特征标之和必须等于一。类似地，$K$ 应当是 $R$ 中性的。矢量多重态是实的，没有自然的 $R$ 对称性。这一对称性在超对称规范理论的研究中起重要作用，下一节我们还会回来讨论它。

% [Tutor Note: 原文 PDF 第 31 页（印刷页 30），$R$ 对称性。注意 $R$ 作用为 $\Phi'(x,\theta)=e^{2in\alpha}\Phi(x,e^{-i\alpha}\theta)$，故等价地 $\theta\to e^{+i\alpha}\theta$、$d^2\theta\to e^{-2i\alpha}d^2\theta$；于是 $\int d^2\theta\,{\cal W}$ 不变要求 ${\cal W}$ 中每一项的 $R$ 特征标之和为 $1$，而 $K$ 须为 $R$ 中性。原文此处 "Since $\theta\to e^{+i\alpha}\theta$, or $d^2\theta\to\ldots$, The R-character..." 中 "or" 应为 "i.e."、"The" 误作大写，均系原文笔误。术语上，本译本把 R-character 统一译作“$R$ 特征标”（即超场的 $R$ 荷）。$R$ 特征标的定义（$A\to e^{2in\alpha}A$，$\psi\to e^{2i(n-1/2)\alpha}\psi$，$F\to e^{2i(n-1)\alpha}F$）与 $d^2\theta$ 的权重 $e^{-2i\alpha}$ 必须配套，否则可重整性对 ${\cal W}$ 的约束会算错。]

\noindent\underline{矢量多重态的拉格朗日量}：
\noindent 如上一节所述，由 (\ref{W-abelian}) 给出的阿贝尔场强 $W$ 是手征超场。利用那里的展开，容易算出
$$
W^\alpha W_\alpha\mid_{\theta\theta}=-2i\lambda\sigma^\mu\partial_\mu
\bar{\lambda} + D^2 - \frac{1}{2} F^{\mu\nu} F_{\mu\nu} +\frac{i}{4} \,
\epsilon^{\mu\nu\rho\sigma} F_{\mu\nu}F_{\rho\sigma}\,.
$$
因此，通常的阿贝尔超对称拉格朗日量（不含 $F\widetilde F$ 项）为
$$
{\cal L}=\frac{1}{4g^{2}}\left(\int d^2\theta\, W^\alpha W_\alpha +
\int d^2\bar\theta\, \bar{W}_{\dot\alpha} \bar{W}^{\dot\alpha}\right)\,.
$$
类似地，在非阿贝尔情形，采用归一化
${\rm Tr}T^aT^b=\delta^{ab}$，有
\begin{equation}
{\rm Tr}(W^\alpha W_\alpha \mid_{\theta\theta})=-2i\lambda^a\sigma^\mu 
D_\mu \bar{\lambda}^a +D^a D^a -\frac{1}{2} F^{a\mu\nu}
F^a_{\mu\nu} +\frac{i}{4}\epsilon^{\mu\nu\rho\sigma} F^a_{\mu\nu} 
F^a_{\rho\sigma}\,,
\label{WW-nonabelian}
\end{equation}
从而通常的非阿贝尔超对称拉格朗日量（不含 $F\widetilde F$ 项）为
$$
{\cal L}=\frac{1}{4g^2}{\rm Tr}\left(\int d^2\theta\, W^\alpha W_\alpha 
+\int d^2\bar\theta\, \bar{W}_{\dot\alpha}\bar{W}^{\dot\alpha}\right)\,.
$$
然而我们关心的是拉格朗日量 (\ref{LXXVIII}) 的超对称类似物，它同时含有 $\theta$ 项。由 (\ref{WW-nonabelian}) 显然可见，带 $\theta$ 项的超杨--米尔斯拉格朗日量可写成
\begin{eqnarray}
{\cal L}&=&\frac{1}{8\pi}{\rm Im}\left(\,\tau\,{\rm Tr}\int d^2\theta\,
W^\alpha W_\alpha \right) \nonumber\\
&=&-\frac{1}{4g^2} F^a_{\mu\nu}F^{a\mu\nu}+ \frac{\theta}{32\pi^2}  
F^a_{\mu\nu}\tilde F^{a\mu\nu}+ \frac{1}{g^2}(\frac{1}{2}D^a D^a 
-i\lambda^a\sigma^\mu D_\mu\bar{\lambda}^a) \,,
\label{tauWW}
\end{eqnarray}
其中 $\tau = \theta/2\pi + 4\pi i/g^2$。注意 $\tau$ 可看作一个常手征超场。

% [Tutor Note: 原文 PDF 第 31 页（印刷页 30），$W^\alpha W_\alpha|_{\theta\theta}$ 与式 (66)--(67)。关键约定：$\widetilde F^{\mu\nu}\equiv\frac12\epsilon^{\mu\nu\rho\sigma}F_{\rho\sigma}$，故式中的 $\frac i4\epsilon^{\mu\nu\rho\sigma}F_{\mu\nu}F_{\rho\sigma}=\frac i2 F\widetilde F$；对偶项的归一化取决于此定义，若把 $\widetilde F$ 定义成无 $1/2$ 的形式，$F\widetilde F$ 的系数会差因子 2。$F\widetilde F$ 为全导数，在通常的超对称拉格朗日量中被丢弃，把它保留并配以 $\tau$ 的虚部（$\theta/2\pi$）即得式 (67)：其中 $-1/4g^2$ 来自 $\tau$ 实部 $4\pi/g^2$ 乘 $-\frac12$ 再乘 $1/8\pi$，$\theta/32\pi^2$ 来自 $\theta/2\pi$ 乘 $\frac12$ 再乘 $1/8\pi$，可用来校验系数。]

\noindent\underline{相互作用项与一般 $N=1$ 拉格朗日量}：
设手征超场 $\Phi_i$ 属于规范群的某个给定表示，其生成元为矩阵 $T^a_{ij}$。动能项
$\Phi_i^\dagger\Phi_i$ 在整体规范变换
$\Phi^\prime=e^{-i\Lambda} \Phi$ 下不变。在局域情形，为保证 $\Phi^\prime$ 仍是手征超场，$\Lambda$ 必须是手征超场。于是超对称且规范不变的动能项为 $\Phi^\dagger e^{-2V}\Phi$。现在我们可以写出完整的 $N=1$ 超对称拉格朗日量：
\begin{equation}
{\cal L}=\frac{1}{8\pi}\,{\rm Im}\left(\tau{\rm Tr}\int d^\theta
\,W^\alpha W_\alpha \right)+\int d^2\theta d^2\bar\theta\,\Phi^\dagger 
e^{-2V}\Phi+ \int d^2\theta\,{\cal W}+\int d^2\bar\theta\,\bar{\cal W}\,.     
\label{fullSF}
\end{equation}
注意，由于每一项各自不变，标量部分与杨--米尔斯部分之间的相对归一化并不由 $N=1$ 超对称固定。上式中我们把标量部分的归一化取为一，但稍后我们将通过重新标度标量多重态 $\Phi$ 来改变它。用分量场表示，上述拉格朗日量取形式
\begin{eqnarray}
{\cal L}\,=&-&\frac{1}{4g^2}F^a_{\mu\nu}F^{a\mu\nu}+\frac{\theta} 
{32\pi^2}F^a_{\mu\nu}\widetilde F^{a\mu\nu}
-\frac{i}{g^2} \,\lambda^a \sigma^\mu D_\mu \bar{\lambda}^a 
+\frac{1}{2g^2}D^a D^a  \nonumber\\
&+& (\partial_\mu A-iA^a_\mu T^a A)^\dagger(\partial^\mu A-iA^{a\mu} 
T^a A) -i\,\bar{\psi} 
\bar{\sigma}^\mu (\partial_\mu\psi -iA^a_\mu T^a \psi)\nonumber\\
&-& D^a A^\dagger T^a A -i\sqrt{2}\,A^\dagger T^a \lambda^a
\psi+i\sqrt{2}\, \bar{\psi} T^a A \bar{\lambda}^a + F^\dagger_i F_i
\nonumber\\ 
&+&\frac{\partial {\cal W}}{\partial A_i} \, F_{i} +
\frac{\partial \bar{\cal W}}{\partial A^\dagger_i} \, F^\dagger_{i} -
\frac{1}{2}\,\frac{\partial^2 {\cal W}}{\partial A_i\partial A_j}\,
\psi_{i}\psi_{j}-\frac{1}{2}\,\frac{\partial^2 \bar{\cal W}} 
{\partial A_i^\dagger\partial
  A_j^\dagger}\,\bar{\psi}_{i}\bar{\psi}_{j}\,. 
\label{fullC}
\end{eqnarray}
这里 ${\cal W}$ 表示超势的标量分量。辅助场 $F$ 与 $D^a$ 可用其运动方程消去；涉及这些场的项从而给出标量势
\begin{equation}
V = \sum_i\,\left|\frac{\partial{\cal W}}{\partial A_i}\right|^2
-\frac{1}{2} g^2 (A^\dagger T^a A)^2\,.
\label{Spot}
\end{equation}

% [Tutor Note: 原文 PDF 第 32 页（印刷页 31），式 (68) 的排版笔误：第一项印作 $\int d^\theta$，缺少上标 $2$，应为 $\int d^2\theta\,W^\alpha W_\alpha$。正文公式按原文照录，正确形式在此注明；若照原文排版，$\int d^\theta$ 只是一个形式上缺了上标的写法，物理含义仍是 $\int d^2\theta$。]

% [Tutor Note: 原文 PDF 第 32 页（印刷页 31），式 (68) 下方。$N=1$ 超对称允许标量部分与杨--米尔斯部分取任意相对归一化（每项分别不变），原文把标量部分归一化为 $1$；下一节（原文 §2.6，见原文 PDF 第 32--33 页式 (69) 之后）将用 $\Phi\to\Phi/g$ 重新标度，使两个费米子的动能项归一化相同，这正是 $N=2$ 对 $N=1$ 拉格朗日量施加的第一个约束。这是理解 $N=2$ 构造的关键一步。]

% [Tutor Note: 原文 PDF 第 32 页（印刷页 31），式 (70) 的符号需核对。由式 (69) 的辅助场项 $+\frac{1}{2g^2}D^aD^a-D^aA^\dagger T^aA$ 消去 $D^a$（得 $D^a=g^2A^\dagger T^aA$）后，得到的应是 $+\frac12g^2(A^\dagger T^aA)^2$：原文式 (70) 第二项写作负号，会使标量势沿 $D$ 方向无下界，且与下一节式 (71) 的正定结果 $-\frac{1}{2g^2}{\rm Tr}([A^\dagger,A]^2)\geq0$ 不一致，疑为原文笔误。译文照录原文的负号，请教授核对原文 PDF 第 32 页式 (70)。]
